\section{Related Work}
\label{sec:related}

\para{Chain-of-thought prompting.} \citet{wei2022chain} introduced CoT prompting, showing that few-shot exemplars with intermediate reasoning steps dramatically improve performance on arithmetic, commonsense, and symbolic reasoning tasks. Their ablations established that the semantic content of the chain matters---replacing reasoning with filler tokens degrades performance---but did not examine whether CoT affects output diversity across runs or models. Zero-shot CoT \citep{kojima2022large} simplified the approach by prepending ``Let's think step by step'' without exemplars. We use zero-shot CoT as our treatment condition because it is the most widely adopted variant and introduces no exemplar-specific bias.

\para{Self-consistency and answer convergence.} \citet{wang2022self} proposed self-consistency, sampling multiple CoT paths and selecting the majority answer. Their key insight---that diverse reasoning paths converge on correct answers while incorrect paths scatter---is the closest prior observation to our findings. However, self-consistency studies convergence \emph{within} CoT (given that CoT is used, do sampled paths agree?) rather than the effect of \emph{adding} CoT (does CoT increase agreement relative to direct prompting?). We address this complementary question. \citet{liu2025answer} extended this line by showing that answers stabilize well before the full CoT completes, with models converging to their final answer after roughly 60\% of reasoning steps. They also observed that extended reasoning can cause ``overthinking'' on easy tasks, consistent with our finding that CoT decreases convergence on multiple-choice tasks where models can answer correctly without explicit reasoning.

\para{Reasoning structure and convergence.} \citet{xiong2025mapping} converted CoT traces into directed reasoning graphs and defined a convergence ratio measuring how reasoning threads are synthesized. They found that convergence ratio strongly predicts accuracy ($r = 0.68$) and that zero-shot prompting produces more convergent reasoning graphs than few-shot prompting. Their structural analysis complements our output-level measurements: while they study the topology of reasoning \emph{within} a single CoT trace, we measure agreement \emph{across} traces and models.

\para{Decoupling of answers and reasoning.} A growing body of work shows that answer-level convergence can coexist with reasoning-level divergence. The MATCHA framework \citep{matcha2025} demonstrated that LLMs maintain correct answers under perturbation while generating inconsistent reasoning chains---the ``Decoupling Hypothesis.'' \citet{chen2023two} formalized two failures of self-consistency: models cannot reliably predict their own outputs (hypothetical consistency) and fail to preserve answers when intermediate results are substituted back in (compositional consistency). These findings suggest that the convergence we measure at the answer level may not reflect convergence in underlying reasoning processes.

\para{Theoretical perspectives on CoT.} \citet{merrill2024representational} proved formally that CoT expands the representational capacity of language models, making them equivalent to probabilistic Turing machines. This expanded capacity could in principle increase output diversity rather than convergence, making our empirical finding---that CoT's convergence effect is task-dependent---all the more notable. \citet{pfau2024dots} showed that even meaningless filler tokens can help transformers solve hard tasks, suggesting that part of CoT's benefit comes from additional computational steps rather than semantic reasoning. This raises the question of whether CoT-induced convergence reflects shared reasoning or shared computation.

\para{Cross-model and cross-context consistency.} \citet{elazar2021measuring} established that pre-trained language models are inherently inconsistent across paraphrased queries, with even the best models achieving only 61\% consistency on factual knowledge probes. \citet{pfau2023self} measured self-consistency on ambiguous tasks, finding that consistency improves with model capability (67--82\% across GPT-3 to GPT-4) but models remain poorly calibrated about their own consistency. Unlike these studies, which measure consistency within a single model, we compare convergence \emph{across} models and examine how prompting strategy modulates it.

\para{Positioning.} Our work differs from all prior studies in a specific way: we treat CoT as an \emph{independent variable} and measure its causal effect on output convergence. Prior work either studies convergence within CoT \citep{wang2022self, xiong2025mapping}, measures consistency without manipulating prompting strategy \citep{elazar2021measuring, pfau2023self}, or analyzes reasoning quality without comparing to a no-CoT baseline \citep{chen2023two, matcha2025}. We bridge this gap by directly comparing the same models on the same questions with and without CoT.
